% Source: https://artofproblemsolving.com/community/c3h2319745_number_theoretic_problem
% Stable identity: aops-topic-2319745-q01
\textbf{(a)} For $2x^2+3x+1=0$ the discriminant is $9-8=1=1^2$, a perfect square, but the roots are $-1$ and $-\dfrac12$, and $-\dfrac12$ is not an integer. So the claim is false. (What is true for integer coefficients is that the roots are \emph{rational} if and only if the discriminant is a perfect square.)

\textbf{(b)} \emph{If the roots are integers.} Let the roots be the integers $u$ and $v$. Then $x^2+bx+c=(x-u)(x-v)$, so $b=-(u+v)$ and $c=uv$, and
\[b^2-4c=(u+v)^2-4uv=(u-v)^2,\]
a perfect square.

\emph{If $b^2-4c$ is a perfect square.} Let $b^2-4c=m^2$ with $m\ge0$ an integer. The roots are
\[x=\dfrac{-b\pm m}{2}.\]
Since $m^2=b^2-4c$, the numbers $m^2$ and $b^2$ differ by a multiple of $4$, so they have the same parity, and hence $m$ and $b$ have the same parity. Then $-b+m$ and $-b-m$ are even, and both roots are integers. $\blacksquare$
